\section{Additional Results and Details}
\label{sec:appendix}

\subsection{Student selection}
\label{app:student-selection}
The student is chosen by successive halving over ten candidate architectures that
share the teacher's tokenizer, so that the cached teacher outputs can be reused
without re-tokenisation. All stages use the validation split only; the test split
is never touched. Stage 0 checks that a candidate fits one 16\,GB GPU and
generates text; stage 1 gives all ten the same short fine-tune; stage 2 retrains
the best five with the full student recipe; stage 3 retrains the best three with
the exact \expid{E2} recipe and two seeds. One candidate is kept. A candidate is
excluded at any stage if it fails one of five gates: it fits the hardware, it uses
the image, it can be compressed (units can be enumerated and removed), its licence
permits the study, and it has seen no test data.
\todo{candidate table, per-stage scores and the selected student, filled from
\texttt{results/} once the selection runs have finished}

% TODO: per-finding sensitivity/specificity for every arm (E0-E9).
% TODO: E9 (teacher 4-bit, matched memory) results.
% TODO: A4 calibration-size curve (32/64/128/256).
% TODO: calibration plots and risk-coverage curves.
% TODO: hyper-parameters (training, recovery fingerprint, PTQ settings).
% TODO: pre-registration document and deviations.
% TODO: compute used (GPU hours per block).
\todo{appendix}
